\section{Eksperymenty i analiza wyników WIP}
\label{sec:Eksperymenty}

\subsection{Metodologia porównania}

Celem eksperymentów jest ilościowe i jakościowe porównanie trzech podejść do zarządzania cyklem życia modelu uczenia maszynowego \cite{kreuzberger2023mlops}:

\begin{enumerate}
  \item \textbf{Podejście ręczne} --- wszystkie kroki (przygotowanie danych, trening, walidacja, wdrożenie) wykonywane przez człowieka bez automatyzacji.
  \item \textbf{Uproszczone CI/CD} --- pipeline z częściową automatyzacją (automatyczny trening po merge, ręczna walidacja jakości i wdrożenie).
  \item \textbf{Zautomatyzowane CI/CD (niniejszy projekt)} --- pełna automatyzacja od push danych do wdrożenia modelu.
\end{enumerate}

Porównanie obejmuje następujące wymiary:
\begin{itemize}
  \item czas potrzebny do wykonania każdego scenariusza,
  \item jakość modelu (metryki Dice i IoU),
  \item pracochłonność (liczba ręcznych kroków),
  \item ryzyko błędu ludzkiego,
  \item skalowalność procesu.
\end{itemize}

\subsection{Czasy wykonania scenariuszy}

Tabela \ref{tab:czasy} przedstawia zmierzone czasy dla każdego podejścia w scenariuszu dodania nowej partii danych treningowych i wdrożenia ulepszonego modelu.

\begin{table}[H]
\centering
\caption{Porównanie czasów wykonania scenariuszy}
\label{tab:czasy}
\begin{tabular}{|l|r|r|r|}
\hline
\textbf{Etap} & \textbf{Ręczny} & \textbf{Uproszczone CI/CD} & \textbf{Zautomatyzowane CI/CD} \\
\hline
Przygotowanie danych & \textit{[uzup.]} & \textit{[uzup.]} & $<1$ min \\
Walidacja danych & \textit{[uzup.]} & \textit{[uzup.]} & \textit{[uzup.]} \\
Trening modelu & \textit{[uzup.]} & \textit{[uzup.]} & \textit{[uzup.]} \\
Ocena jakości & \textit{[uzup.]} & \textit{[uzup.]} & $<1$ min \\
Wdrożenie & \textit{[uzup.]} & \textit{[uzup.]} & \textit{[uzup.]} \\
\hline
\textbf{Łącznie} & \textit{[uzup.]} & \textit{[uzup.]} & \textit{[uzup.]} \\
\hline
\end{tabular}
\end{table}

% [UZUPEŁNIĆ: komentarz do wyników tabeli]

\subsection{Metryki modelu}

Tabela \ref{tab:metryki} przedstawia metryki modelu dla kolejnych wersji wytrenowanych w ramach eksperymentów. Baseline pochodzi z repozytorium referencyjnego \textit{Water-Meters-Segmentation} \cite{kaggle_water_meters}. Metryki są rejestrowane automatycznie przez MLflow \cite{zaharia2018mlflow} podczas każdego treningu.

\begin{table}[H]
\centering
\caption{Metryki jakości modelu dla kolejnych wersji}
\label{tab:metryki}
\begin{tabular}{|l|c|c|l|}
\hline
\textbf{Wersja modelu} & \textbf{Dice} & \textbf{IoU} & \textbf{Uwagi} \\
\hline
Baseline (referencyjny) & 0{,}9275 & 0{,}8865 & Oryginalny model \\
v1 --- pierwsze trenowanie & \textit{[uzup.]} & \textit{[uzup.]} & \textit{[uzup.]} \\
v2 --- po dodaniu danych & \textit{[uzup.]} & \textit{[uzup.]} & \textit{[uzup.]} \\
\hline
\end{tabular}
\end{table}

% [UZUPEŁNIĆ: screenshoty z MLflow UI z wykresami metryk]

\subsection{Metryki infrastruktury}

\subsubsection{Zużycie zasobów podczas treningu}

% [UZUPEŁNIĆ: screenshoty z Grafana - CPU/RAM podczas treningu na EC2]

Podczas treningu modelu monitorowano obciążenie instancji EC2. Kluczowe obserwacje:
\begin{itemize}
  \item szczytowe zużycie CPU: % [UZUPEŁNIĆ] \%,
  \item szczytowe zużycie RAM: % [UZUPEŁNIĆ] GB z % [UZUPEŁNIĆ] GB dostępnych,
  \item czas treningu jednej epoki: % [UZUPEŁNIĆ] s.
\end{itemize}

\subsubsection{Zużycie zasobów podczas serwowania}

% [UZUPEŁNIĆ: screenshoty z Grafana - metryki API w trakcie testów obciążeniowych]

API modelu serwującego predykcje charakteryzuje się następującymi parametrami wydajnościowymi:
\begin{itemize}
  \item średni czas odpowiedzi dla jednego obrazu: % [UZUPEŁNIĆ] ms,
  \item przepustowość: % [UZUPEŁNIĆ] żądań/sekundę,
  \item zużycie RAM przez pod: % [UZUPEŁNIĆ] MB.
\end{itemize}

\subsubsection{Koszty infrastruktury AWS}

\begin{table}[H]
\centering
\caption{Szacowane koszty infrastruktury AWS}
\label{tab:koszty}
\begin{tabular}{|l|r|l|}
\hline
\textbf{Zasób} & \textbf{Koszt} & \textbf{Uwagi} \\
\hline
EC2 t3.large (trening) & \textit{[uzup.]} USD/h & Uruchamiana tylko podczas treningu \\
EC2 t3.large (serwowanie) & \textit{[uzup.]} USD/h & Stop/start przez workflow \\
S3 (artefakty) & \textit{[uzup.]} USD/mies. & Zależy od liczby modeli \\
ECR (obrazy) & \textit{[uzup.]} USD/mies. & \\
\hline
\textbf{Łącznie (projekt)} & \textit{[uzup.]} USD & Za cały okres realizacji \\
\hline
\end{tabular}
\end{table}

\subsection{Ocena jakościowa}

\subsubsection{Łatwość wdrożenia}

% [UZUPEŁNIĆ: ocena na podstawie doświadczeń z implementacji]

\subsubsection{Stabilność systemu}

W trakcie testów zaobserwowano następujące zachowania systemu w sytuacjach wyjątkowych:
\begin{itemize}
  \item \textbf{Niepoprawne dane} --- walidator odrzucał dane z komunikatem i nie uruchamiał treningu,
  \item \textbf{Model gorszy od baseline} --- bramka jakości odrzucała model bez tworzenia PR,
  \item \textbf{Awaria instancji EC2} --- job \texttt{stop-infra} uruchamiał się zawsze (\texttt{if: always()}), zapobiegając wyciekom kosztów.
\end{itemize}

\subsubsection{Skalowalność}

Zastosowane podejście ma następujące ograniczenia skalowalności:
\begin{itemize}
  \item single-node k3s --- brak wysokiej dostępności,
  \item jeden runner EC2 --- brak równoległych treningów,
  \item model proof of concept --- nie testowano przy zbiorach $>1000$ obrazów.
\end{itemize}
